\scanpage{94}
a component $D_\alpha=\overline{\{s_\alpha\}}$ of this divisor, is nothing other than its \uline{multiplicity} along $D_\alpha$.

This recalled, one therefore wonders about \unsure{that|in} which happens for the $s\in\bigcup D_\alpha$, \deleted{\unsure{One|(the}} \inserted{and} \deleted{sees that} \inserted{\unsure{to look at.|let us look at}} those which are $\in D_\alpha$, $\alpha$ \unsure{taken|fixed|fixed}, i.e. which are specializations of a fixed $s_\alpha$. One can say that the situation defines an extension $\mathcal E_{U_{s_\alpha}}$ of $\mathcal E_{k(s_\alpha)}$ by $\pi$ --- deduced, as \unsure{earlier|earlier} for $S$ and $K$, from a \unsure{relative|relative} curve $U_{D_\alpha}$ over $D_\alpha$ --- or $\unsure{U_{S_\alpha}|U_{s_\alpha}}$ over a suitable regular open subset of $D_\alpha$, if one prefers --- by passing to the generic point. This time, for the generic fiber, one no longer has \deleted{\unsure{sections|section}} rational points over $k(s_\alpha)$ of $\unsure{U_{s_\alpha}|U_{S_\alpha}}$ --- one has only one point: the point at infinity $i_\alpha\in X_{s_\alpha}\setminus U_{s_\alpha}$, rational over $\unsure{k_\alpha\xi^{\cdot}|K_\alpha\xi^{\cdot}}$. But one \inserted{now} has a relative $L_i\subset\pi$: this $i_\alpha$, precisely, in an essentially canonical way,
\margin{93}
